\paragraph{Detecting frames on the admitted cover.}
Choose the coefficient embedding used to construct the cover.
A nonzero algebraic section has nonzero pullback to at least one component of the complex uniformization. The marked period domains cover the complex moduli
components, and its values in frames determine the section on each component. On a component where it does not vanish identically, let \(s(y)\) be its normalized modulus as a function of the frame. Choose \(y_0\) there with
\[
 s(y_0)>\tfrac12\sup_y s(y)>0.
\]
Cuspidality bounds the section. If \(H\) is the positive majorant at this frame and \(L_{\rm ref}\) the reference lattice, discreteness and positive definiteness give 
\[
 \min_{0\ne v\in L_{\rm ref}}H(v,v)>0.
\]

The positive minimum gives the lattice lower bound for the interior CM approximation used in
\hyperref[subsec:deep-cm-interior]{the lower-norm argument}.
Let \(u_1,\ldots,u_a\) be the simultaneous approximants to the positive axes
in the Hermitian space \((V,h)\) over \(E\). The two spaces used in the construction are
\[
 P=\operatorname{span}_E\{u_1,\ldots,u_a\},\qquad
 P^{\perp_h}=\{v\in V:h(v,u_j)=0\text{ for }1\le j\le a\}.
\]
The short-disc argument preserves the nonzero value of the section. Its input is the
period lattice in the uniformization of the deep cover, related to
the reference lattice by the fractional ideal in the frame
construction. The resulting adapted basis supplies the two CM
types, their Riemann lattice and the original marked polarization.

The CM realization uses this lattice and its level-\(N\) marking.
After a finite coefficient extension \(K/F\) it gives a CM point
\(z\) with an integral CM model and with the specified complex base
point. The comparison uses the marked fibre isomorphism with the
universal abelian variety, so the complex point belongs to the
same component of the very cover from which we started.
\label{cor:deep-cover-denominators-1}

\begin{lemma}[Algebraic frame detecting a section]
\label{lem:deep-cover-detecting-frames}
For a section \(s\) on an admitted component, let \(\mathcal F\) denote the algebraic interior CM frames preserving its marking. Then
\[
 s\ne0\quad\Longrightarrow\quad \exists\,y\in\mathcal F,\qquad s(y)\ne0.
\]

The detecting frame is allowed to depend on the section.
\end{lemma}

\begin{proof}
\proofstep{Moving a detected frame to an algebraic point.}
The given nonzero complex frame may be different from the
algebraic reference frame \(y_0\) of \(z\). Fix the theta generator
first. On the variety of block frame changes over \(K\), evaluation
is a regular function of the matrix entries with the block
determinants inverted. A change of the theta generator multiplies
it by a nonzero scalar given by the central exponent. Consequently
the rational function in the matrix entries is nonzero, since
its complex evaluation is nonzero.

If \(N\) is the numerator and \(Y_\tau\) the block variables, consider
\[
 P(Y)=N(Y)\prod_\tau\det Y_\tau\ne0.
\]
A nonzero value of \(P\) tests the numerator and all invertibility conditions simultaneously. A polynomial over an infinite field cannot vanish on
all tuples of its elements, by induction on the number of
variables. It follows that some tuple over \(K\) gives an
invertible block change with a nonzero evaluation. This is the
detecting frame \(y\). The argument applies also at a nonzero
central exponent, which only alters the theta factor of the
frame change.
\label{cor:deep-cover-denominators-2}

\end{proof}

\paragraph{One denominator for the evaluation at a frame.}

A rational frame detects a nonzero section. Passing from that evaluation to an integral test requires a denominator that accounts for every graded-line exponent.

Write \(e_c\) for the central exponent and \(e_{\tau i}\) for the
flag exponents, with \(m_\tau\) letters in block \(\tau\). For fixed
\(z,y_0\) and exponents put
\[
 k_\tau=\sum_{i=1}^{m_\tau}\max(0,-e_{\tau i}),\qquad
 D_e=\sum_{\tau,i}(e_{\tau i}+k_\tau).
\]
For every section \(f\) of the weight there is a polynomial
\(P_f\) in the block variables, of total degree \(D_e\) when
nonzero, such that
\[
 P_f(Y)=\prod_\tau\det(Y_\tau)^{k_\tau}\,f(Y\cdot y_0).
\]
We use the convention that \(Y\) is the transpose of the
frame-change matrix, so the formula fixes the coefficient
convention. Here is the reason the indicated determinant factors
suffice. The evaluation is regular on the invertible blocks and
diagonal scaling gives column weights \(e_{\tau i}\) in the
transposed convention. A uniform determinant twist adds
\(k_\tau\) to each of these weights. In a block, a possible
remaining denominator is a power of its irreducible
determinant.

On the open subset of its singular divisor where
the first \(m_\tau-1\) columns are independent, right
multiplication by an upper unitriangular matrix subtracts a
linear combination of those columns from the last. The
resulting last column is a local parameter times a fixed
transverse column on an open chart. Weight equivariance under
scaling that column gives order
\[
 e_{\tau,m_\tau}+k_\tau\ge0
\]
along the divisor.
The other blocks may be taken invertible on this chart.
The argument applies separately to each block
and also to a one-dimensional block. In the polynomial ring
there is therefore no remaining determinant denominator.
Simultaneous scaling of all block entries now gives degree
\(D_e\). In particular only finitely many monomials can occur,
independently of \(f\).
\label{cor:deep-cover-denominators-3}

\begin{lemma}[One denominator at a fixed frame]
\label{lem:deep-cover-frame-clearing}
For a fixed section \(s\) and algebraic frame \(y\), let \(b=s(y)\) be its evaluation in the indicated coordinates. There is one nonzero coefficient scalar \(c\) such that
\[
 v(cb)\ge0\qquad\text{at all finite valuations }v.
\]
 The comparison uses its homogeneous block
polynomial and the determinant factors appearing in the frame
change.
\end{lemma}

\begin{proof}
\proofstep{The local check of the frame polynomial.}
At a finite place \(w\) of \(K\), complete the integral CM
model. It has an integral eigenbasis for the individual
eigenspace lattices and a unit generator for the local theta
line. These data yield the valued frame tests for the same
cover, also after an extension of the valued field. In this
frame the determinant-twisted polynomial of a section with
nonnegative test orders has integral coefficients. Indeed, for
integral frame-change matrices of unit block determinants the
evaluations are integral.

If \(P\) had a coefficient \(a\) of negative minimal order, normalization would give
\[
 v(a)<0,\qquad a^{-1}P\in\mathcal O_w[Y],\qquad \overline{a^{-1}P}\ne0.
\]
The reduced polynomial times the block determinants is
nonzero. It has a nonzero evaluation after passing to a large
enough residue extension, contradicting the order inequality in
the lifted integral frames. This uses the full family of valued
tests. An integral basis here is taken in each individual
eigenspace, which can be compared with the reference letters
even at a place where the decomposition into eigenspaces is
only a generic-fibre decomposition.
\label{cor:deep-cover-denominators-4}

\proofstep{Returning to the original frame.}
To return to \(y_0\), let \(G_\tau\) be the transition blocks
to the integral eigenbasis and let \(s\ne0\) compare the theta
generators, in the direction where the new generator is \(s\)
times the old. Choose nonzero constants \(c_{\tau j}\) clearing
the rows of \(G_\tau^{-1}\) when returning the reference letters
to that basis. Substitution into the integral polynomial and
the theta and determinant transformations show that
\[
 \left(\prod_\tau\det(G_\tau)^{k_\tau}\right)^{-1}
 s^{e_c}\left(\prod_{\tau,i,j}c_{\tau j}^{\mu_{\tau ij}}\right)
 [Y^\mu]P_f\quad\text{is integral at }w .
\]
This is a coefficient statement. Each column rescaling
contributes once for each occurrence of an entry of that column
in the monomial. Substitution of the remaining integral
matrices introduces no additional denominator.

For any exceptional \(w\), the finitely many possible \(\mu\)
give a single local lower bound on all coefficient orders.
Outside the places where the fractional letter and theta ideals
or the determinant comparison have support, the letter scalars
are \(1\) and \(s\) is a unit. In this case the inverse
transition blocks are integral and the product of their
original block determinants is integral. For each \(\tau\) the required individual determinant is
\[
 \det G_\tau=\Bigl(\prod_\sigma\det G_\sigma\Bigr)\prod_{\sigma\ne\tau}\det(G_\sigma^{-1}).
\]
Every factor on the right is integral. The displayed
formula then gives integral coefficients at every such \(w\).
Thus the finite exceptional set is fixed by the CM data and
the weight before varying \(f\). At each exceptional place
take a power of a nonzero integer in its prime ideal of
sufficient order. Their product \(t\) clears the coefficients
simultaneously at all places.
\label{cor:deep-cover-denominators-5}

For the algebraic detecting frame \(y\) let \(Y\) be the
transition tuple from \(y_0\) and \(c\) the central scalar.
Then
\[
 f(y)=c^{e_c}\prod_\tau\det(Y_\tau)^{-k_\tau}P_f(Y).
\]
The values of the finitely many degree-\(D_e\) monomials
at \(Y\), times the fixed multiplier in this expression,
have a common denominator after multiplication by \(t^{-1}\).
It gives the required scalar for all integral \(f\). The
scalar may be taken to be a positive rational integer
by clearing in \(K\). This proof also works for the shifted
weights since \(e_c\) and the exponents were fixed integers
throughout.
\label{cor:deep-cover-denominators-6}

\end{proof}

\paragraph{Matching a large nonnegative weight.}
Let \(B=4[E:\Q]r+4\). The cubical invariant on the marked
cover has a fourth root \(\rho\), with \(\rho^4=a_Z\mathcal I\)
on each component \(Z\), where \(a_Z\in F^\times\) and
\(\mathcal I\) has nonnegative test orders. Multiply by a
single \(t_1\in\mathcal O_F\setminus\{0\}\) with
\[
 t_1^4a_Z\in\mathcal O_F
\]
for every \(Z\).
The fourth-power identity gives nonnegative orders for
\(t_1\rho\). Thus \(u=\rho^{-1}\) and \(u'=t_1\rho\) have
product \(t_1\), with \(u'\) integral at all tests.
The exponents of \(u^B\) add central column count \(2B\)
and subtract \(B\) in every flag coordinate.

At every label a parallel flat determinant unit and its
inverse are bounded on the complex component, and both
have a finite clearing scalar there. Take the common
positive period \(a\) of their parallel exponent steps.
Their products and inverses realize a prescribed shift
\(aw_\tau\) at each label, in the opposite directions on
the two blocks in the flag exponents. Clearing the
inverse on all finitely many components gives a pair
\(\phi,\phi'\) with \(\phi\phi'=t_2\in F^\times\) and
\(\phi'\) integral at every test.
\label{cor:deep-cover-denominators-7}

\begin{lemma}[Clearing a high-column form]
\label{lem:deep-weight-shifts}
Suppose the two block weights of the cusp form \(f\) have
gap \(\lambda_p-\lambda_{p+1}\ge1\) at each label. Let \(B\) be
the common determinant shift and \(a>0\) the period of the parallel
flat units. For every integer \(b\ge2B+a\) one can choose \(w_\tau\in\mathbb Z\) so that, for \(F_b=f^bu^B\phi\), each block weight \(\lambda_{\tau,i}(F_b)\) and the column count \(n(F_b)\) satisfy
\[
 \min_{\tau,i}\lambda_{\tau,i}(F_b)\ge0,\qquad n(F_b)=2B.
\]
 The determinant and flat
multipliers are generically invertible on the component. The two complementary short
expansions then give one clearing scalar for the original
high-column form at all valuations, with no uniform bound on
its size as the form or level varies.
\end{lemma}

\begin{proof}
\proofstep{The possible determinant shifts.}
At a label write the nonincreasing block weight as
\(\lambda_1,\ldots,\lambda_{p+q}\) with
\(\lambda_p-\lambda_{p+1}\ge1\).
Choose an integer \(b\ge 2B+a\) and write
\[
 b\lambda_p-B=am+R,\qquad 0\le R<a.
\]
Use \(w_\tau=-m\) here. In the \(p\)-block the target
exponents \(b\lambda_i+aw_\tau-B\) are at least \(R\).
In the conjugate block, in reverse index order, the target
exponents \(-b\lambda_i-aw_\tau-B\) are at least
\[
 b-2B-R\ge0.
\]
In both blocks they form nonincreasing
sequences. This gives an admissible nonnegative weight of even
column count \(2B\) for \(f^b u^B\phi\). Its inverse multiplier
up to scalar is \((u')^B\phi'\).

\proofstep{Preservation of the cusp condition.}
To check the cusp condition, factor off a single \(f\).
If the remaining section is nonzero, its block exponents
are dominant by the vanishing of sections with a negative
adjacent gap on a flag fibre. Its central exponent is even.
The framed uniformization and the fourth-root comparison
therefore make it holomorphic with the corresponding
weight. With \([E:\Q]/2>2\), Koecher's argument on the
unitary period domains makes its Fourier--Jacobi support
nonnegative at each cusp. Write the product as \(fg\), where \(f\) is the
original cusp factor and \(g\) is the remaining section.
Fourier--Jacobi expansion gives
\[
    [fg]_0=\sum_{\nu+\mu=0}[f]_\nu[g]_\mu
           =[f]_0[g]_0=0 .
\]
The pointed nonnegative support gives the middle equality
and the cusp condition on \(f\) gives the last. Thus the product
belongs to the high-column cusp space. The zero product
satisfies the same assertion.
\label{cor:deep-cover-denominators-8}

\paragraph{One scalar for all valuations of a high-column form.}
We use the short theta generator comparison on the odd marked
level \(N\), with auxiliary level and comparison modulus both
one. The added comparison rule is the constant rule one, and
its torsion coefficient has prescribed trivial character. The
uniformized components have the same rational Hermitian form
and the same marked polarization vector. Their principal
congruence actions may differ by central roots of unity. Let \(w_E=|\mu(E)|\) and replace the high-column form \(h\) by
\(h^{w_E}\). Every resulting frame exponent \(a\), including the
half-column contribution, is divisible by \(w_E\), so the discrepancy
in the central action disappears:
\[
                     \zeta^a=1\qquad(\zeta\in\mu(E)).
\] The column count only increases.

\proofstep{Two theta expansions with different prime conditions.}
Here are the two comparisons used for this form. For any
positive integer \(b_0\), the short theta forms with an
auxiliary unimodular lower-row lattice trivial at the primes
of \(b_0N\) span its complex realization. The hypotheses on the column count \(n\) include 
\[
 n\in2\mathbb Z,\qquad n\ge4[E:\Q]r+4.
\]
The target weights are dominant. Raising to the power \(w_E\)
above kills the central multiplier, and the additional target
character is trivial. To impose the auxiliary conditions while
keeping this target, set
\[
 B=2b_0N,\qquad S=\mathcal O_E^n,\qquad
 H=U(S\otimes_{\mathcal O_E}E)
\]
for the standard positive definite Hermitian form.
The definite unitary integral in
Proposition~\ref{prop:span-original-theta} uses a source marking
at \(B\). Its doubled kernel is in the convergent range of the
Siegel--Weil formula since \(n\ge4[E:\mathbb Q]r+4\).
\cite[Theorem~2]{yamana-siegel-weil}

For a finite adelic representative \(h=(h_v)_v\) the lattice of
the class has local components \(h_vS_v\). At every \(v\mid B\)
prescribe the following coset for the new representative:
\[
 K_v(B)=\{k\in H(E^+_v):kS_v=S_v,\ (k-1)S_v\subset BS_v\},
 \qquad u_v\in K_v(B)h_v^{-1}.
\]
To choose a single rational \(u\), use the Cayley chart
\[
 X^*=-X,\qquad C(X)=(1-X)(1+X)^{-1}.
\]
Its complement in the group, given by \(\det(1+u)=0\), has empty
interior at each finite place. Choose a parameter in each prescribed
open coset and approximate the parameters in the skew-Hermitian
vector space over \(E^+\). The resulting \(u\in H(E^+)\) gives
a lattice \(\Lambda\) with
\[
 \Lambda_v=S_v\quad(v\mid B),\qquad \Lambda^\vee=\Lambda .
\]
Let \(e_1,e_2\) be the exponents of the finite groups
\(S/(S\cap\Lambda)\) and \(\Lambda/(S\cap\Lambda)\). Their
localizations vanish at the primes of \(B\). Hence with
\(e=\operatorname{lcm}(e_1,e_2)\) we have
\[
 (e,B)=1,\qquad e\Lambda\subset S,\qquad eS\subset\Lambda .
\]

The rational unitary change reindexes the summation in the theta
kernel. Its real action preserves harmonicity and the target row
weight. Refining a compact open source level \(K\) to \(K'\)
preserves the integral of a function invariant under \(K\), since
\[
           \sum_{k\in K/K'}\operatorname{vol}(kK')
                       =[K:K']\operatorname{vol}(K')
                       =\operatorname{vol}(K).
\]
Use these representatives in the inverse-overlap and reproducing
calculation of Proposition~\ref{prop:span-original-theta}, before
the good-type projection. The target half-level is fixed. Extend
the component functions to their distinct orbits as at
\ref{cor:slice-shifted-theta-1}. The extra target ideal is the unit
ideal, so the lower-row congruence is the same for the arithmetic
and analytic evaluations. Expanding each polynomial in integral
harmonic rows gives the required short rules on the same components.

Let \(W\) be the \(F\)-span of the rules in the section space \(V\). Scalar extension gives the test
\[
 v\in W\quad\Longleftrightarrow\quad (v+W)\otimes1=0\text{ in }(V/W)\otimes_F\mathbb C.
\]
Indeed an
\(F\)-basis of the subspace generated by the rules can be
extended to a basis of the section space, and any nonzero
coordinate outside the subspace would remain nonzero at the
coefficient embedding.
\label{cor:deep-cover-denominators-9}

For each resulting short rule \(g\), choose its nonzero scalar
\(t_g\) and positive auxiliary integer \(n_g\). The two sets of
primes at which the scalar clears all test denominators are
\[
       \{\ell:\ell\mid b_0\}\ \cup\ \{\ell:\ell\nmid n_g\}.
\]
The same \(t_g\) works throughout this union. Here \(n_g\) is the
odd integer for its theta sum. To see the scope of the
integrality calculation, let \(\ell\) be a prime at which the
unimodular lattice is trivial and write
\(N=\ell^cN'\), \(n_g=\ell^{c+d}n'\), with
\(\ell\nmid N'n'\). The \(\ell\)-primary part of the torsion
classes is read using \(\ell^c\) and the prime-to-\(\ell\)
modulus \(n'\).

The replacement isomorphism on the local
lattice uses \(\ell^d\) and its Hermitian norm
\(\ell^{2d}\), removing the extra depth in the index and
polynomial denominators of the short sum. The class estimate
then costs only the level power, the discriminant factor of the CM field
and the fixed generator comparison. In the notation of the
short sum, with flag degree \(D\), these are cleared by
\[
 t_g=\beta_H(\Lambda)\,
       \#\{\text{level tuples}\}\,
       N^{D}\prod_{\ell\mid C_E}\ell^{|\operatorname{disc}E|D}.
\]
The normalized local order inequality holds on the lifted
integral frame, and normalized orders are preserved by the
valued base change. The trivial-lattice condition holds at
\(\ell\mid b_0\) by construction and at \(\ell\nmid n_g\)
by the lattice's level condition.
\label{cor:deep-cover-denominators-10}

\proofstep{A clearing scalar on the remaining tests.}
Now use a finite expansion with \(b_0=1\). Replacing its
generators by \(t_g g\) and clearing the coefficients in
\(F\) gives a scalar \(s_1\in\mathcal O_F\setminus\{0\}\)
making the powered form integral away from the primes
dividing
\[
 n_* =\prod_g n_g.
\]
Use a second finite expansion
with \(b_0=n_*\). It gives a scalar
\(s_2\in\mathcal O_F\setminus\{0\}\) making the same
powered form integral at all tests over the primes dividing
\(n_*\). Every test thus gives nonnegative order after
multiplication by \(s_1s_2\), since the other factor is
integral.

If the power was \(w\ge1\), then whenever the values have finite valuation,
\[
 wv(s_1s_2 h)=v(s_1s_2 h^w)+(w-1)v(s_1s_2)\ge0 .
\]
A zero value is integral as well. This provides the
simultaneous denominator for the original high-column
form, using a possibly different short expansion for the
exceptional prime set. None of these finite denominators
needs a uniform size as \(N\), the weight or the form varies.
\label{cor:deep-cover-denominators-11}
\end{proof}
